% ============================================================
% COMMANDS.TEX
% Sistema Adaptativo de Classificação Multicritério
% de Fluidos de Corte
%
% FINALIDADE
%
% Centralizar comandos reutilizáveis da apresentação.
%
% Este arquivo faz parte apenas da BASE do projeto.
%
% Não inserir aqui:
%
% - resultados;
% - pesos CRITIC;
% - coeficientes TOPSIS;
% - ranking;
% - conclusões;
% - valores calculados.
%
% Resultados científicos deverão ser produzidos posteriormente
% pelos integrantes por meio da execução em R.
%
% ============================================================


% ============================================================
% PLACEHOLDER PADRÃO
% ============================================================
%
% Utilizado enquanto uma informação ainda não foi confirmada.
%
% Exemplo:
%
% \TBF
%
% ============================================================

\providecommand{\TBF}{\texttt{TO\_BE\_FILLED}}


% ============================================================
% IDENTIFICAÇÃO DO PROJETO
% ============================================================

\providecommand{\ProjectTitle}{
Sistema Adaptativo de Classificação Multicritério de Fluidos de Corte
}

\providecommand{\ProjectSubtitle}{
Auditoria de Dados, CRITIC, TOPSIS, Pareto e Análise de Robustez
}

\providecommand{\ProjectYear}{
2026
}


% ============================================================
% INFORMAÇÕES INSTITUCIONAIS
% ============================================================
%
% Preencher somente quando as informações oficiais estiverem
% confirmadas pela professora ou pelo grupo.
%
% ============================================================

\providecommand{\InstitutionName}{
\TBF
}

\providecommand{\CourseName}{
\TBF
}

\providecommand{\DepartmentName}{
\TBF
}

\providecommand{\DisciplineName}{
\TBF
}

\providecommand{\ProfessorName}{
\TBF
}


% ============================================================
% INTEGRANTES
% ============================================================
%
% Nesta fase utilizamos somente os nomes pelos quais os
% integrantes estão sendo identificados no projeto.
%
% Antes da apresentação final, substituir pelos nomes completos
% quando necessário.
%
% ============================================================

\providecommand{\MemberPacheco}{
Pacheco
}

\providecommand{\MemberBenjamin}{
Benjamin
}

\providecommand{\MemberVitor}{
Vitor
}

\providecommand{\MemberMarco}{
Marco
}


% ============================================================
% LISTA DE INTEGRANTES
% ============================================================

\providecommand{\ProjectMembers}{
\MemberPacheco,
\MemberBenjamin,
\MemberVitor,
\MemberMarco
}


% ============================================================
% NOMES DOS MÉTODOS
% ============================================================
%
% Estes comandos servem apenas para manter a nomenclatura
% consistente ao longo da apresentação.
%
% ============================================================

\providecommand{\CRITIC}{
\textsc{CRITIC}
}

\providecommand{\TOPSIS}{
\textsc{TOPSIS}
}

\providecommand{\PCA}{
\textsc{PCA}
}


% ============================================================
% TERMOS PADRONIZADOS
% ============================================================

\providecommand{\BenefitCriterion}{
\texttt{BENEFIT}
}

\providecommand{\CostCriterion}{
\texttt{COST}
}

\providecommand{\TargetCriterion}{
\texttt{TARGET}
}

\providecommand{\RestrictionCriterion}{
\texttt{RESTRICTION}
}

\providecommand{\DiagnosticCriterion}{
\texttt{DIAGNOSTIC}
}


% ============================================================
% STATUS DO PROJETO
% ============================================================

\providecommand{\StatusDraft}{
\texttt{DRAFT}
}

\providecommand{\StatusReview}{
\texttt{UNDER\_REVIEW}
}

\providecommand{\StatusApproved}{
\texttt{APPROVED}
}

\providecommand{\StatusBlocked}{
\texttt{BLOCKED}
}

\providecommand{\StatusSuperseded}{
\texttt{SUPERSEDED}
}


% ============================================================
% STATUS DE EXECUÇÃO
% ============================================================

\providecommand{\NotExecuted}{
\texttt{NOT\_EXECUTED}
}

\providecommand{\NotGenerated}{
\texttt{NOT\_GENERATED}
}

\providecommand{\NotDefined}{
\texttt{NOT\_DEFINED}
}


% ============================================================
% TERMOS DE CONFORMIDADE
% ============================================================

\providecommand{\Confirmed}{
\texttt{CONFIRMED}
}

\providecommand{\PendingConfirmation}{
\texttt{PENDING\_CONFIRMATION}
}

\providecommand{\UnknownStatus}{
\texttt{UNKNOWN}
}


% ============================================================
% TIPOS DE RESTRIÇÃO
% ============================================================

\providecommand{\HardConstraint}{
\texttt{HARD\_CONSTRAINT}
}

\providecommand{\SoftConstraint}{
\texttt{SOFT\_CONSTRAINT}
}


% ============================================================
% TEXTO DE ALERTA PARA RESULTADOS AINDA NÃO GERADOS
% ============================================================
%
% Pode ser utilizado durante a construção da apresentação.
%
% Exemplo:
%
% \PendingResult
%
% Deve ser removido ou substituído na versão final apresentada.
%
% ============================================================

\providecommand{\PendingResult}{
\textit{Resultado a ser preenchido após a execução da análise.}
}


% ============================================================
% TEXTO DE ALERTA PARA MATERIAL EXTERNO
% ============================================================

\providecommand{\ExternalMaterialPending}{
\textit{Material externo a ser inserido após confirmação.}
}


% ============================================================
% MACROS MATEMÁTICAS — CRITIC
% ============================================================
%
% Apenas atalhos de notação.
%
% Não executam cálculo.
%
% ============================================================

\providecommand{\CriticInformation}[1]{
C_{#1}
}

\providecommand{\CriticWeight}[1]{
w_{#1}
}

\providecommand{\CriticSD}[1]{
\sigma_{#1}
}

\providecommand{\CriterionCorrelation}[2]{
r_{#1#2}
}


% ============================================================
% MACROS MATEMÁTICAS — TOPSIS
% ============================================================

\providecommand{\TopsisPositiveDistance}[1]{
S_{#1}^{+}
}

\providecommand{\TopsisNegativeDistance}[1]{
S_{#1}^{-}
}

\providecommand{\TopsisCoefficient}[1]{
C_{#1}
}


% ============================================================
% MACRO PARA SETAS DE PIPELINE
% ============================================================

\providecommand{\PipelineArrow}{
\(\rightarrow\)
}


% ============================================================
% MACRO PARA TEXTO DE REGRA
% ============================================================
%
% Uso:
%
% \ProjectRule{texto}
%
% Mantém a apresentação visual simples, sem depender de cores
% personalizadas definidas neste arquivo.
%
% ============================================================

\providecommand{\ProjectRule}[1]{
\begin{alertblock}{Regra}
#1
\end{alertblock}
}


% ============================================================
% MACRO PARA PRINCÍPIO
% ============================================================

\providecommand{\ProjectPrinciple}[1]{
\begin{block}{Princípio}
#1
\end{block}
}


% ============================================================
% MACRO PARA RESULTADO FUTURO
% ============================================================
%
% Uso:
%
% \FutureResult{Nome do resultado}
%
% ============================================================

\providecommand{\FutureResult}[1]{
\begin{block}{#1}
\TBF
\end{block}
}


% ============================================================
% RESULTADOS — NÃO DEFINIR AQUI
% ============================================================
%
% NÃO criar agora comandos como:
%
% \WinnerFluid
% \TopOne
% \TopTwo
% \TopThree
% \HighestCriticWeight
% \TopsisScoreA
%
% porque esses valores ainda não existem oficialmente.
%
% Se futuramente o grupo decidir automatizar a inserção dos
% resultados no LaTeX, esses comandos deverão ser gerados a
% partir dos resultados do pipeline em R, e não preenchidos
% manualmente.
%
% ============================================================


% ============================================================
% CAMINHOS
% ============================================================
%
% Não colocar caminhos absolutos como:
%
% C:/Users/Nome/Desktop/...
%
% Utilizar sempre caminhos relativos ao projeto.
%
% ============================================================


% ============================================================
% LOGOS
% ============================================================
%
% Os arquivos oficiais de logo deverão ficar, conforme a
% estrutura atual da apresentação, em:
%
% templates/presentation/latex/assets/
%
% Os nomes dos arquivos somente deverão ser definidos depois que
% os materiais efetivos forem adicionados.
%
% Portanto não criamos macros de caminho de logo agora.
%
% ============================================================


% ============================================================
% COMPATIBILIDADE
% ============================================================
%
% Os comandos deste arquivo utilizam \providecommand em vez de
% \newcommand para reduzir conflitos caso algum comando seja
% definido anteriormente pelo template ou tema.
%
% ============================================================


% ============================================================
% RESPONSABILIDADE DOS INTEGRANTES
% ============================================================
%
% Durante a execução do projeto, os integrantes poderão:
%
% - acrescentar macros realmente úteis;
% - automatizar a inserção de valores;
% - simplificar comandos;
% - remover placeholders;
%
% desde que preservem:
%
% - rastreabilidade;
% - consistência;
% - reprodutibilidade.
%
% ============================================================


% ============================================================
% ESTADO ATUAL
% ============================================================
%
% COMMANDS_STATUS = BASE_PREPARED
%
% SCIENTIFIC_RESULTS_DEFINED = FALSE
% INSTITUTION_DATA_CONFIRMED = FALSE
% OFFICIAL_LOGOS_DEFINED = FALSE
%
% ============================================================